\documentclass[../../main]{subfiles}

\begin{document}
\section{Regras de Inferência}
\label{sec:matematica:logica-matematica:regras-de-inferencia}

\begin{defn}[Regra de Inferência]
	\label{def:matematica:logica-matematica:regras-de-inferencia:regra-de-inferencia}

	Uma regra de inferência é um esquema lógico que permite obter uma
	conclusão a partir de um conjunto de premissas de modo que a verdade
	das premissas implique necessariamente a verdade da conclusão.
\end{defn}

\begin{defn}[Modus Ponens]
	\label{def:matematica:logica-matematica:regras-de-inferencia:modus-ponens}

	Sejam \(P\) e \(Q\) proposições.

	O esquema

	\[
		P,
	\]

	\[
		P \Rightarrow Q
	\]

	implica

	\[
		Q.
	\]

\end{defn}

\begin{prop}[Validade do Modus Ponens]
	\label{prop:matematica:logica-matematica:regras-de-inferencia:validade-modus-ponens}

	O Modus Ponens é uma regra de inferência válida.
\end{prop}

\begin{proof}

	Considere a implicação

	\[
		\bigl(
		P
		\land
			(P \Rightarrow Q)
		\bigr)
		\Rightarrow
		Q.
	\]

	Sua tabela-verdade possui valor verdadeiro para toda atribuição de
	valores de verdade.

	Logo, trata-se de uma tautologia.

	Portanto, o Modus Ponens é válido.

\end{proof}

\begin{defn}[Modus Tollens]
	\label{def:matematica:logica-matematica:regras-de-inferencia:modus-tollens}

	Sejam \(P\) e \(Q\) proposições.

	O esquema

	\[
		P \Rightarrow Q,
	\]

	\[
		\neg Q
	\]

	implica

	\[
		\neg P.
	\]

\end{defn}

\begin{prop}[Validade do Modus Tollens]
	\label{prop:matematica:logica-matematica:regras-de-inferencia:validade-modus-tollens}

	O Modus Tollens é uma regra de inferência válida.
\end{prop}

\begin{proof}

	Considere

	\[
		\bigl(
		(P \Rightarrow Q)
		\land
		\neg Q
		\bigr)
		\Rightarrow
		\neg P.
	\]

	Verifica-se por tabela-verdade que a expressão é uma tautologia.

	Portanto, o Modus Tollens é válido.

\end{proof}

\begin{defn}[Silogismo Hipotético]
	\label{def:matematica:logica-matematica:regras-de-inferencia:silogismo-hipotetico}

	Sejam \(P\), \(Q\) e \(R\) proposições.

	O esquema

	\[
		P \Rightarrow Q,
	\]

	\[
		Q \Rightarrow R
	\]

	implica

	\[
		P \Rightarrow R.
	\]

\end{defn}

\begin{prop}[Validade do Silogismo Hipotético]
	\label{prop:matematica:logica-matematica:regras-de-inferencia:validade-silogismo-hipotetico}

	O Silogismo Hipotético é uma regra de inferência válida.
\end{prop}

\begin{proof}

	Considere

	\[
		\Bigl(
		(P \Rightarrow Q)
		\land
		(Q \Rightarrow R)
		\Bigr)
		\Rightarrow
		(P \Rightarrow R).
	\]

	A expressão é uma tautologia.

	Portanto, o Silogismo Hipotético é válido.

\end{proof}

\begin{defn}[Silogismo Disjuntivo]
	\label{def:matematica:logica-matematica:regras-de-inferencia:silogismo-disjuntivo}

	Sejam \(P\) e \(Q\) proposições.

	O esquema

	\[
		P \lor Q,
	\]

	\[
		\neg P
	\]

	implica

	\[
		Q.
	\]

\end{defn}

\begin{prop}[Validade do Silogismo Disjuntivo]
	\label{prop:matematica:logica-matematica:regras-de-inferencia:validade-silogismo-disjuntivo}

	O Silogismo Disjuntivo é uma regra de inferência válida.
\end{prop}

\begin{proof}

	Considere

	\[
		\bigl(
		(P \lor Q)
		\land
		\neg P
		\bigr)
		\Rightarrow
		Q.
	\]

	A expressão é uma tautologia.

	Portanto, o Silogismo Disjuntivo é válido.

\end{proof}

\begin{defn}[Redução ao Absurdo]
	\label{def:matematica:logica-matematica:regras-de-inferencia:reducao-ao-absurdo}

	Seja \(P\) uma proposição.

	Se a hipótese \(\neg P\) conduz a uma contradição, então conclui-se
	que \(P\) é verdadeira.

	Simbolicamente,

	\[
		\neg P
		\Rightarrow
		\bot
	\]

	implica

	\[
		P.
	\]

\end{defn}

\begin{prop}[Validade da Redução ao Absurdo]
	\label{prop:matematica:logica-matematica:regras-de-inferencia:validade-reducao-ao-absurdo}

	A Redução ao Absurdo é uma regra de inferência válida na lógica
	clássica.
\end{prop}

\begin{proof}

	Pela Lei da Dupla Negação e pelo Princípio do Terceiro Excluído,
	demonstrados anteriormente, segue que

	\[
		(\neg P \Rightarrow \bot)
		\Rightarrow
		P.
	\]

	Portanto, a regra é válida na lógica clássica.

\end{proof}
\end{document}
